\chapter{Financiación Tradicional}
\section{Introducción}
\paragraph{¿Cuándo una empresa necesita financiación?}
Para poder realizar las operaciones normales de la empresa (las NOF: necesidades operativas de fondos). Normalmente se financia parte con Fondo de Maniobra y parte con crédito o deuda a corto.

Las empresas suelen necesitar financiación a largo plazo (Deuda a largo plazo o capital) cuando se enfrentan a los siguientes desafíos:
\begin{itemize}
    \item Iniciar un proyecto empresarial: obtener los recursos iniciales para desarrollar e implementar una idea de negocio (etapa altamente riesgosa).
    \item Falta de liquidez para mantener un proyecto empresarial: obtener los fondos necesarios para afrontar los pagos (riesgo de insolvencia, situación financiera desfavorable, poca o nula rentabilidad).
    \item Crecer como empresa: Obtener los recursos para invertir en los diferentes proyectos y así expandir el negocio (situación económica financiera favorable).
\end{itemize}

\paragraph{Mapa de fuentes de financiación}
\begin{itemize}
    \item Internas
    \begin{itemize}
        \item Resultados (positivos) de las Operaciones (Flujos de caja / Cash Flow)
        \item Capital Social (Ampliaciones)
    \end{itemize}
    \item Externas
    \begin{itemize}
        \item Pasivos espontáneos (proveedores, otras ctas a pagar)
        \item Corto plazo
        \begin{itemize}
            \item Línea de crédito (crédito comercial)
            \item Factoring
            \item Confirming
        \end{itemize}
        \item Largo plazo
        \begin{itemize}
            \item Préstamo bancario
            \item Leasing, Renting
            \item Subvenciones
        \end{itemize}
    \end{itemize}
\end{itemize}

\section{Fuentes de Financiación Bancarias}
Las fuentes privadas de financiación se pueden agrupar en dos categorías: a largo plazo y a corto plazo.

\subsection{Corto Plazo}
\subsubsection{Créditos Comerciales}
También llamadas líneas o cuentas de crédito son un tipo de crédito bancario que da ifnanciamiento a corto plazo a una empresa.

La empresa dispondrá de un dinero en su cuenta y accederá a él cuando le sea necesario. A medida que vaya utilizándolo irá agotando dicho crédito, teniendo un límite. Periódicamente la empresa deberá de devolver dicho crédito a corto plazo. Normalmente la empresa pagará intereses a medida que vaya utilizando el dinero.

normalmente incluyen intereses deudores, excendidos y acreedores. Y varias comisiones.

\subsubsection{Factoring}
Es un sistema de financiación enfocado principalmente para empresas que tienen que gestionar cobros de facturaws. Por lo que la empresa anticipa el cobro de sus pagos pendientes y da por finalizada la transacción comercial con el cliente.

Permite a una empresa acreedora disponer del importe anticipadamente a la fecha de vencimiento del cobro.

Tras intervinientes, la empresa de Factoring, el cliente y el deudor.

La empresa que ofrece los servicios de Factoring puede asumir el riesgo comercial. Si lo asume se denomina Factoring sin recurso.

La gestión de cobro abarca la fase de impagados.

De gran utilidad para cobros por exportaciones.

\subsubsection{Confirming}
Una empresa contrae una deuda con un proveedor y tiene que proceder al pago. Dicho pago deberá ser efectuado a 30 días. Pero, para esa fecha, la empresa cree que no va a disponer de liquidez suficiente. Por lo tanto, acude al confirming. Una entidad financiera soportará el pago de la empresa, a cambio de un interés fijado. Los costes son mayores cuanto menos solvente sea tu empresa ante el banco.

\subsubsection{Préstamos}
El préstamo bancario es el tipo de financiación más utilizado y el más conocido. Este sistema va muy ligado a la solvencia financiera que tenga la empresa y su propio patrimonio personal, Es decir, las garantías que pueda aportar al banco toman mucha importancia.
\begin{itemize}
    \item Se concede un determinado importe (principal).
    \item Se establece un plan para devolver el principal y abonar los intereses (vencimientos).
    \item Se determina el tipo de interés nominal para la liquidación.
    \item Puede tener varias comisiones: de apertura, de cancelación anticipada, de amortización anticipada.
\end{itemize}
Puede llevar asociado un serguro de reembolso por fallecimiento, invalidez o desempleo.

\subsection{Largo Plazo}
\subsubsection{Leasing (Arrendamiento financiero)}
Es un contrato mercantil de arrendamiento con opción de compra.

A la fecha de vencimiento se puede ejrcitar la opción de compra, prorogar el contrato o solicitar la reposición del bien por otro de tecnología más avanzada.

Permite a una empresa disponer de activos fijos que son propiedad de la Entidad Financiera.

Se establece a medio o largo plazo.

La Entidad Financiera los adquiere para la finalidad establecid apor el arrendatario.

Permite financiar el 100\% del valor del activo.

Leasing = Alquiler + Opción de Compra

\paragraph{Fiscalidad Leasing}
\begin{itemize}
    \item La carga financiera satisfecha (intereses) es gasto fiscalmente deducible.
    \item La parte satisfecha para la recuperación del coste del bien también es fiscalmente deducible salvo activos no amortizables, con un límite anual.
    \item La deducción es al tipo del impuesto de sociedades.
\end{itemize}

\subsubsection{Renting (Arrendamiento operativo)}
El renting es un sistema de financiación en el que se alquila un bien, y pasado un tiempo establecido, se puede devolver y cambiar por uno nuevo. El renting no está regulado por la Ley. Sin embargo, sí que tiene también, beneficios fiscales para la empresa. En el caso del renting, la propiedad del bien no se transmite.

El renting es un mecanismo de contratación para adquirir servicios, maquinaria o productos con los que generalmente no cuenta una empresa. Otra forma de decirlo: el renting es una herramienta que tienen las empresas para adquirir productos que necesitan como si de un alquiler/arrendamiento se tratara y que por ahora no quieren o no pueden permitirse su adquisición (compra).

\begin{itemize}
    \item Es un alquiler operativo a medio y largo plazo para bienes de equipo, vehículos u otros bienes empresariales a cambio de una renta periódica.
    \item Empresa arrendadora (especializada en renting) entrega el bien, lo gestiona, mantiene y asegura (servicio integral).
    \item Empresa arrendataria disfruta del bien y al finalizar el plazo lo devuelve.
    \item En general para bienes con obsolescencia rápida, con mantenimiento y que no forman parte de los procesos básicos de la emrpesa.
    \item Se debe cumplir el plazo establecido. La cancelación anticipada tiene fuertes penalizaciones.
    \item La empresa de Renting asume un doble riesgo del cliente y del bien.
\end{itemize}

\paragraph{Diferencias con el Leasing}
\begin{itemize}
    \item El Renting lo pueden utilizar empresas y particulares.
    \item En el Renting se suele devolver el bien al finalizar el contrato (aunque se puede pacta opción de compra) y en el Leasing se puede comprar, renovar o devolver.
    \item El grado de utilización del bien determina el importe de la renta (vehículos y kilometraje), loq eu no sucede en el Leasing.
    \item Las cuotas del Renting suelen ser más altas que las del Leasing debido a los servicios añadidos de mantenimiento.
    \item Contablemente el Renting se considera un gasto (no requiere recursos propios para la asquisición) y el Leasing afecta a las cuentas de pasivo por el importe de la deuda yy al activo como gastos diferidos.
\end{itemize}

Las fuentes de financiación externas son todo aquellos fondos que ofertan los diversos actores de los mercados financieros y que se conceden a cambio d eque la empresa genera valor. Pueden ser de origen público o privado y a largo o corto plazo.
\begin{itemize}
    \item Créditos ICO: son préstamos ofrecidos por el Ministerio de Enocomía, cuyo principal objetivo es poder financiar a los autónomos, pequeñas sociedades y entidades públicas y privadas. Máximo 10 millones de euros.
    \item CDTI: Es la entidad que canaliza las solicitudes de ayuda y apoyo a los proyectos de I+D+i de empresas españolas en los ámbitos estatal e internacional.
    \item ENISA: Dirección General de Industria y de la Pequeña y Mediana Empresa, integrada, a su vez, en el Ministerio de Industris, Comercio y Turismo.
    \item Subvenciones: la Administración Pública ofrece ayudas o bien por sectores (argicultura, industria, turismo...) o por autonomías.
    \item INEM: solicitar el pago único de la prestación del desempleo para financiar tu propio negocio. 
\end{itemize}

\section{Fuentes de Financiación no Bancarias}
\subsection{Mercados}

\subsection{Deuda}
\subsubsection{Financiación interna}
Es la financiación que obtieen la empresa con los propios recursos generados mediante la actividad que realiza. El capital social de una empresa puede estar constituido por diferentes tipos de acciones:
\begin{itemize}
    \item Acciones ordinarias: representan la propiedad de una francción del capital.
    \item Acciones liberadas: acciones ordinarias que se entragan sin desembolso del suscriptor o con un desembolso parcial.
    \item Acciones privilegiadas: otorgan derechos o ventajas (mayor porcentaje de dividendos, mayor votos en la junta de accionistas).
    \item Acciones sin voto: privan al suscriptor del derecho a voto a cambio de otros privilegios.
    \item Acciones rescatables: otorgan el derecho a la empresa de rescatar esta accionistas después de un tiempo.
\end{itemize}

\subsubsection{Valoración de acciones}
\[Capital\ social = N * n\]
Valor de las acciones:
\begin{itemize}
    \item Valor nominal: valor convencional otorgado al título en el momento de la emisión, si bien puede modificarse a lo largo del tiempo.
    \item Valor contable (o teórico-contable): es el resultado de dividir el Patrimonio Neto entre el número de acciones existentes.
    \item Valor razonable: es el precio que se paga por la acción en el mercado.
\end{itemize}

Ampliación de capital:
\begin{itemize}
    \item Emitiendo más acciones ($N$) y/o
    \item Aumentando el valor nominal de las existentes ($n$)
\end{itemize}

Reducción de capital:
\begin{itemize}
    \item Devolución de aportaciones de los socios
    \item Condonación de dividendos pasivos/deuda de los socios
    \item Constitución o incremento de reservas/permuta
    \item Saneamiento de pérdidas
\end{itemize}


% En una empresa, si se amplia el capital aumentando el valor nominal de las acciones existentes (n), ¿los accionistas deben pagar la diferencia del valor nominal respecto a lo que pagaron?

% Si mis acciones las he sacado a 10€/acción y quiero ampliar el capital